\chapter{Introduction to Digitisation of Historical Inscriptions and Manuscripts}
This chapter frames the digitisation challenge for South Asian historical inscriptions and manuscripts and situates the project within heritage preservation and modern AI-assisted workflows. It introduces the project title and context, then moves into the motivation, problem statement, and objectives that drive the work. A brief methodology overview outlines the staged pipeline from preprocessing and AI enhancement to binarisation. The section on assumptions and constraints clarifies the operational boundaries, and the chapter closes by summarizing how the rest of the report is organized.

\section[Introduction]{\textbf{Introduction}}
South Asia preserves a vast body of historical knowledge in stone inscriptions, palm leaf manuscripts, copper plates, paper manuscripts, and cave or rock paintings. Many of these artefacts are fragile, geographically dispersed, and visually degraded, which makes the scanned images available in digital archives difficult to study directly. While major archives provide scanned images, the quality of those images is often insufficient for legible reading or downstream transcription: degradation by noise, uneven illumination, fading, and surface weathering renders large portions of each scan illegible.

This project, titled \textit{Digitisation of Historical Inscriptions and Manuscripts}, addresses that gap through a unified, multi-artefact image processing pipeline. The workflow takes degraded scans and produces an enhanced image and a clean binary output that separates text pixels from the background, providing a visually legible and machine-processable representation of each artefact. Unlike single-script or single-medium systems, the design routes each artefact type through tailored enhancement and binarisation steps to preserve text legibility and material context.

The work presented in this report covers Phase 1 of the pipeline: preprocessing, AI enhancement, and binarisation. Downstream processing stages are planned for Phase 2 and beyond, after the image quality foundation established here has been validated. This section sets the technical and historical context for the project and anchors the motivation, problem statement, and objectives that follow.

\section[Literature Review]{\textbf{Literature Review}}
This review focuses on recent work in historical document enhancement, binarisation, character recognition, layout analysis, and translation, with an emphasis on studies most relevant to South Asian inscriptions and manuscripts.

Aitha et al. present a Brahmi character dataset and a custom CNN for Ashokan Brahmi, which is a critical step toward character recognition for ancient scripts but remains an early-stage dataset effort without integration into a full processing pipeline \cite{Aitha2023}.

Mehrotra et al. introduce Indiscapes for layout analysis of Indic manuscripts using CNN-based segmentation, yet the scope is limited to page layout and does not address enhancement, binarisation, or transcription quality \cite{Mehrotra2022}.

AI4Bharat's Bhashini initiative provides transformer-based NLP resources for Indian languages, but its focus is modern language usage and it provides limited support for classical or epigraphic forms \cite{AI4Bharat2022}.

Wang et al. propose Real-ESRGAN for blind super-resolution on real degraded photos, demonstrating strong perceptual restoration, but the evaluation is not conducted on inscription imagery or document-specific degradations \cite{Wang2021}.

Pramanik et al. apply U-Net based denoising to degraded handwritten documents, achieving clear noise suppression, yet the experiments are centered on European scripts rather than Indic historical sources \cite{Pramanik2021}.

Valy et al. build a deep learning pipeline for Khmer palm leaf manuscripts using CNN and LSTM for character recognition, but the approach is language-specific and is not validated on South Asian scripts or inscription media \cite{Valy2021}.

Seuret et al. propose the DIVA-HisDB framework for historical document analysis with layout and transcription components, though the focus is European medieval manuscripts without Indic script coverage \cite{Seuret2021}.

Dhanya et al. release the THPLMD dataset and apply Sauvola adaptive thresholding for Tamil palm leaf binarisation, which improves foreground extraction but stops short of transcription or translation stages \cite{Dhanya2020}.

Helsinki-NLP's OPUS-MT models enable transformer-based translation for Dravidian and Indic languages, yet performance degrades on archaic or classical text varieties common in inscriptions \cite{HelsinkiNLP2020}.

Krishnan et al. use deep CNNs for character recognition of Tamil and Devanagari, providing script-specific recognition but not incorporating ensemble strategies or confidence-driven fusion \cite{Krishnan2019}.

Burie et al. report the ICFHR 2016 competition on palm leaf manuscripts, benchmarking binarisation methods and highlighting the difficulty of noise and texture, but without any AI enhancement or end-to-end transcription \cite{Burie2016}.

Harman introduces DStretch for rock art enhancement through color decorrelation, a highly effective visual method that is not paired with character recognition or structured text extraction \cite{Harman2008}.

Collectively, these studies show strong individual advances in enhancement, binarisation, character recognition, layout analysis, and translation, but also reveal a gap in integrated, multi-artefact image processing pipelines tuned to South Asian epigraphy and manuscript collections. The present project specifically targets this image-quality gap, delivering a validated preprocessing, enhancement, and binarisation foundation on which downstream stages can be built in Phase 2.

\section[Motivation]{\textbf{Motivation}}
Preserving inscriptions and manuscripts is essential for historical scholarship, yet the materials are often degraded, fragile, and dispersed across archives, limiting access and slowing research. Existing digital collections typically provide only raw scanned images, and the visual quality of those images is frequently insufficient for legible reading or reliable downstream processing. Recent advances in AI-based image enhancement make it feasible to recover faint character detail, suppress noise, and correct uneven illumination automatically. This project is motivated by the need for an end-to-end, repeatable image processing pipeline that improves the visual legibility of degraded artefact scans and produces clean binary images that form a reliable foundation for future transcription and preservation work.

\section[Problem statement]{\textbf{Problem statement}}
Historical inscriptions and manuscripts in South Asia remain largely inaccessible because their scans are visually degraded by noise, uneven illumination, fading ink, surface weathering, and substrate aging. Current archival workflows provide raw images but offer no automated path from degraded scans to enhanced, legible images with clean binary separation of text from background. The problem is to design and implement a reliable, multi-artefact image processing pipeline that can enhance degraded images, binarise text regions using adaptive methods, and validate the quality of the outputs through objective image quality metrics, while respecting non-destructive processing constraints.

\section[Objectives]{\textbf{Objectives}}
The objectives of the project are
\begin{enumerate}
\item To develop a non-destructive, five-artefact preprocessing and AI enhancement pipeline for stone inscriptions, palm leaf manuscripts, copper plates, paper manuscripts, and cave/rock paintings, with artefact-specific processing chains for improved image legibility.
\item To produce high-quality binary images through adaptive Sauvola binarisation with morphological refinement, providing clean text-background separation as a foundation for downstream Phase~2 processing.
\item To validate the quality of enhanced and binarised outputs using objective image quality metrics — PSNR, SSIM, CNR, and edge sharpness — across all five artefact categories.
\end{enumerate}

\section[Brief Methodology of the project]{\textbf{Brief Methodology of the project}}
The project follows a three-stage methodology designed for historical image degradation and artefact diversity. Raw scans are first normalised in preprocessing (orientation correction, brightness normalisation, white balance, and border cleanup). The cleaned images are then enhanced using non-local means denoising, AI super-resolution via Real-ESRGAN, and DStretch-based decorrelation enhancement where material characteristics demand colour separation. Enhanced outputs are converted to binary text-focused images using Sauvola adaptive thresholding, followed by morphological closing and noise blob removal. Image quality is evaluated at the enhancement stage using PSNR, SSIM, CNR, and edge sharpness metrics. The pipeline endpoint is a clean binary image ready for downstream Phase~2 processing.

\section[Assumptions made / Constraints of the project]{\textbf{Assumptions made / Constraints of the project}}
The validity of this project is based on the following assumptions and constraints:
\begin{itemize}
\item Input artefacts are already available as scanned or photographed 2D images in standard formats (JPG/PNG/TIFF), and hardware acquisition is outside project scope.
\item Original archival images are treated as read-only assets; all processing is non-destructive and writes only to derived output directories.
\item Phase 1 implementation covers preprocessing, AI enhancement, and binarisation only. Further stages are deferred to Phase 2 and beyond.
\item Performance observations are based on the selected dataset mix and available compute resources, and may vary with different archives, scan quality, or deployment environments.
\end{itemize}

\section[Organization of the report]{\textbf{Organization of the report}}
This report is organized as follows.
\begin{itemize}
\item Chapter 2 establishes the theoretical foundation for the digitisation pipeline, covering degradation mechanisms across all five artefact categories, image processing fundamentals (colour spaces, CLAHE, Real-ESRGAN super-resolution, DStretch decorrelation stretch, and binarisation theory), ECE signal processing contributions (Gabor filter banks and FFT-based noise removal), and the objective image quality metrics (PSNR, SSIM, CNR, and edge sharpness) used for pipeline evaluation.
\item Chapter 3 presents the architectural design of the three-stage pipeline, including the functional requirements and quality targets, per-stage algorithmic design decisions (preprocessing parameters, Real-ESRGAN configuration, Sauvola binarisation), artefact-specific processing chains, a consolidated design decisions table, and the non-destructive processing policy.
\item Chapter 4 details the stage-by-stage implementation across all Phase~1 modules (\texttt{preprocess.py}, \texttt{enhance.py}, and \texttt{binarise.py}), the ECE signal processing and quality modules, and the 5-test continuous integration suite.
\item Chapter 5 reports quantitative evaluation results across the 38-image test set, presenting PSNR, SSIM, CNR, and edge sharpness tables for all five artefact categories, per-stage throughput profiling, and a qualitative analysis of before/after enhancement outcomes supported by key inferences.
\item Chapter 6 synthesises the Phase~1 outcomes against each project objective with supporting quantitative evidence, discusses the current limitations, outlines the planned Phase~2 extensions (LaMa inpainting and 3D inscription processing), and enumerates the learning outcomes across all three project disciplines.
\end{itemize}
